\documentclass{beamer}
\usefonttheme{serif}
\usepackage{graphicx}
\usepackage{xcolor}
\usepackage{colortbl}
\usepackage{multirow,tabularx}
\usepackage{tikz}
\usepackage{hyperref}
\usepackage{physics}
\usepackage{siunitx}
\usepackage{multicol}
\usepackage{wrapfig}
\usepackage{amsmath}
\usepackage{amsthm}

\hypersetup{
    colorlinks=true,
    linkcolor=blue,
    filecolor=blue,
    urlcolor=blue,
    pdftitle={goodFeaturesToTrack},
    pdfpagemode=FullScreen,
}

% ---------------- COLORS ----------------
\definecolor{blue}{rgb}{0.023, 0.113, 0.609}
\setbeamercolor{frametitle}{fg=blue}
\setbeamercolor{title}{fg=blue}
\setbeamercolor{structure}{fg=blue}

% ---------------- LISTINGS ----------------
\usepackage{listings}

\lstdefinelanguage{python}{
    morekeywords={if, else, for, while, def, return, import, print, in,
                  True, False, None, try, except, class},
    morecomment=[l]{\#},
    morestring=[b]",
    morestring=[b]'
}

\lstset{
    language=python,
    basicstyle=\ttfamily\small,
    keywordstyle=\color{blue}\bfseries,
    commentstyle=\color{gray},
    stringstyle=\color{red},
    numbers=left,
    numberstyle=\tiny\color{gray},
    frame=single,
    breaklines=true,
    backgroundcolor=\color{gray!10}
}

% ---------------- TITLE ----------------
\title{\textcolor{blue}{Feature Detection in Computer Vision}}
\subtitle{\textcolor{blue}{Wigner Summer Camp \\ Data and Compute Intensive Sciences Research Group}}
\author{\textcolor{blue}{Balázs, Paszkál, Vince, Levente, Bence, Antal \\ Éva, Hajni}}
\date{\textcolor{blue}{6--10 June 2026}}

% ---------------- FOOTLINE ----------------
\setbeamertemplate{footline}{
  \leavevmode%
  \hbox{%
  \begin{beamercolorbox}[wd=.1\paperwidth,ht=2.25ex,dp=1ex,left]{page number}%
    \hspace{1em}\textcolor{blue}{\insertframenumber}
  \end{beamercolorbox}}%
  \vskip0pt%
}

% ---------------- DOCUMENT ----------------
\begin{document}

% TITLE
\begin{frame}
    \titlepage
\end{frame}

% ---------------- INTRO ----------------
\begin{frame}{What is feature detection?}
    \begin{itemize}
        \item Detects important points in images
        \item Finds stable structures (corners, edges)
        \item Used for tracking and motion estimation
    \end{itemize}
\end{frame}

% ---------------- METHOD ----------------
\begin{frame}{goodFeaturesToTrack}
    \begin{itemize}
        \item Detects strong corner-like points
        \item Based on intensity changes in both directions
        \item Selects stable points for tracking
    \end{itemize}

    \vspace{0.3cm}

    Output: set of $(x, y)$ coordinates
\end{frame}

% ---------------- THEORY ----------------
\begin{frame}{Corner intuition}
    A good feature has strong variation in both directions:

    \begin{equation}
        \lambda_1, \lambda_2 \gg 0
    \end{equation}

    \begin{itemize}
        \item $\lambda_1, \lambda_2$ = eigenvalues of gradient matrix
        \item Large values → strong corner
        \item Flat regions → not selected
    \end{itemize}
\end{frame}

% ---------------- OPEN CV CODE ----------------
\begin{frame}[fragile]{OpenCV implementation}
    \begin{lstlisting}
import cv2

img = cv2.imread("frame.png")
gray = cv2.cvtColor(img, cv2.COLOR_BGR2GRAY)

corners = cv2.goodFeaturesToTrack(
    gray,
    maxCorners=100,
    qualityLevel=0.01,
    minDistance=10
)

for c in corners:
    x, y = c.ravel()
    cv2.circle(img, (int(x), int(y)), 3, (0,255,0), -1)
\end{lstlisting}
\end{frame}

% ---------------- PARAMETERS ----------------
\begin{frame}{Parameters}
    \begin{itemize}
        \item \textbf{maxCorners} – max number of detected points
        \item \textbf{qualityLevel} – threshold for acceptance
        \item \textbf{minDistance} – minimum spacing between points
        \item \textbf{blockSize} – neighborhood size
    \end{itemize}
\end{frame}

% ---------------- APPLICATIONS ----------------
\begin{frame}{Applications}
    \begin{itemize}
        \item Optical flow tracking (Lucas-Kanade)
        \item Camera stabilization
        \item Robotics perception
        \item Video analysis
    \end{itemize}
\end{frame}



\end{document}